\documentclass[12pt,a4paper]{article}
\usepackage[utf8]{inputenc}
\usepackage[T5]{fontenc}
\usepackage{amsmath, amssymb}
\usepackage{graphicx}
\usepackage{ifthen}

\newboolean{showsolutions}
\setboolean{showsolutions}{true}

% --- ĐỊNH NGHĨA CÁC LỆNH ẢO CHO CÔNG CỤ LATEX2JSON CỦA WEBSITE ---
\newcommand{\doctitle}[1]{\begin{center}\Large\textbf{#1}\end{center}}
\newcommand{\docdesc}[1]{\textbf{Mô tả:} #1}
\newcommand{\docgrade}[1]{\textbf{Lớp:} #1}
\newcommand{\docstatus}[1]{\textbf{Trạng thái:} #1}
\newcommand{\doctype}[1]{\textbf{Loại:} #1}
\newcommand{\doctopics}[1]{\textbf{Chủ đề:} #1}

\newenvironment{textblock}{\vspace{1em}}{}

\newenvironment{lesson}[2]{
	\vspace{1em}\noindent\textbf{\Large #1}\\
	\textit{#2}\vspace{0.5em}
}{}

\newcommand{\image}[3]{
	\begin{center}
		\includegraphics[width=#1\textwidth]{#3} \\
		\textit{#2}
	\end{center}
}

\newenvironment{quiz}[1]{
	\vspace{1em}\noindent\textbf{\Large Kiểm tra: #1}
}{}

\newenvironment{mcq}[4]{
	\vspace{1em}\noindent\textbf{Câu hỏi:} #1 \\
	\ifthenelse{\boolean{showsolutions}}{
		\textit{(Đáp án đúng: #2, Điểm: #3) - Giải thích: #4}
	}{}
	\begin{itemize}
	}{
	\end{itemize}
}
\newcommand{\option}[1]{\item #1}

\newenvironment{truefalse}[3]{
	\vspace{1em}\noindent\textbf{Đúng/Sai:} #1 \\
	\ifthenelse{\boolean{showsolutions}}{
		\textit{(Điểm: #2) - Giải thích: #3}
	}{}
	\begin{itemize}
	}{
	\end{itemize}
}
\newcommand{\statement}[2]{\item [\textbf{#1}] #2}

\newcommand{\shortanswer}[4]{
    \vspace{1em}\noindent\textbf{Trả lời ngắn (Điểm: #3):}

    \noindent #1

	\ifthenelse{\boolean{showsolutions}}{
		\vspace{0.5em}
		\noindent\textit{Đáp án: #2 - Giải thích:}
		\noindent #4
	}{}
}

\newcommand{\essay}[3]{
    \vspace{1em}\noindent\textbf{Tự luận (Điểm: #2):}
    
    \noindent #1
    
	\ifthenelse{\boolean{showsolutions}}{
		\vspace{0.5em}
		\noindent\textbf{Gợi ý / Đáp án:}
		\noindent #3
	}{}
}

\begin{document}

\doctitle{PHẦN III: BÀI TẬP RÈN LUYỆN - VECTƠ TRONG KHÔNG GIAN}
\docdesc{Tài Liệu Ôn Thi Group - FLASH STUDY}
\docgrade{12}
\docstatus{draft}
\doctype{normal}
\doctopics{Hình học không gian, Vectơ trong không gian}

% =========================================================================
% 1. CÂU TRẮC NGHIỆM NHIỀU PHƯƠNG ÁN LỰA CHỌN
% =========================================================================

\begin{quiz}{Câu trắc nghiệm nhiều phương án lựa chọn}

\begin{mcq}{Cho hình hộp $ABCD.A'B'C'D'$. Mệnh đề nào sau đây là đúng?\image{0.6}{}{lt_1.png}}{0}{1}{Ta có $\vec{AC'} + \vec{CA'} = (\vec{AC} + \vec{CC'}) + (\vec{CA} + \vec{AA'}) = 2\vec{CC'} = -2\vec{C'C} \Rightarrow \vec{AC'} + \vec{CA'} + 2\vec{C'C} = \vec{0}$.}
	\option{$\vec{AC'} + \vec{CA'} + 2\vec{C'C} = \vec{0}$.}
	\option{$\vec{CA'} + \vec{AC'} = \vec{CC'}$.}
	\option{$\vec{AC'} + \vec{AC} = 2\vec{AC}$.}
	\option{$\vec{AC'} + \vec{A'C} = \vec{AA'}$.}
\end{mcq}

\begin{mcq}{Cho tứ diện đều $ABCD$. Góc giữa 2 vectơ $\vec{AB}, \vec{BD}$ là:\image{0.6}{}{lt_2.png}}{3}{1}{Ta có $(\vec{AB}, \vec{BD}) = 180^\circ - (\vec{BA}, \vec{BD}) = 180^\circ - 60^\circ = 120^\circ$.}
	\option{$30^\circ$.}
	\option{$60^\circ$.}
	\option{$90^\circ$.}
	\option{$120^\circ$.}
\end{mcq}

\begin{mcq}{Cho hình lập phương $ABCD.A'B'C'D'$. Gọi $G$ là trọng tâm tam giác $A'BC$. Mệnh đề nào sau đây là đúng?\image{0.6}{}{lt_3.png}}{3}{1}{Ta có $\vec{BA'} + \vec{BB} + \vec{BC} = 3\vec{BG} \Rightarrow \vec{BA'} + \vec{BC} = 3\vec{BG}$. Mặt khác $\vec{BD'} = \vec{BA'} + \vec{BC}$, do đó $\vec{BD'} = 3\vec{BG}$.}
	\option{$\vec{AC'} = 3\vec{AG}$.}
	\option{$\vec{AC'} = 4\vec{AG}$.}
	\option{$\vec{BD'} = 4\vec{BG}$.}
	\option{$\vec{BD'} = 3\vec{BG}$.}
\end{mcq}

\begin{mcq}{Cho hình lập phương $ABCD.A'B'C'D'$ có cạnh bằng $a$. Tính độ dài vectơ $\vec{x} = \vec{AA'} + \vec{AC'}$ theo $a$.\image{0.6}{}{lt_4.png}}{2}{1}{Ta có $\vec{x} = \vec{AC} + 2\vec{AA'}$. Vì $\vec{AA'} \perp \vec{AC}$ nên $|\vec{x}|^2 = AC^2 + 4AA'^2 = 2a^2 + 4a^2 = 6a^2 \Rightarrow |\vec{x}| = a\sqrt{6}$.}
	\option{$a\sqrt{2}$.}
	\option{$(1 + \sqrt{3})a$.}
	\option{$a\sqrt{6}$.}
	\option{$\frac{a\sqrt{6}}{2}$.}
\end{mcq}

\begin{mcq}{Cho hình lập phương $ABCD.A'B'C'D'$. Gọi $O$ là tâm hình lập phương. Mệnh đề nào sau đây là đúng?\image{0.6}{}{lt_5.png}}{1}{1}{Vì $O$ là trung điểm đường chéo $AC'$ nên $\vec{AO} = \frac{1}{2}\vec{AC'} = \frac{1}{2}(\vec{AB} + \vec{AD} + \vec{AA'})$.}
	\option{$\vec{AO} = \frac{1}{3}(\vec{AB} + \vec{AD} + \vec{AA'})$.}
	\option{$\vec{AO} = \frac{1}{2}(\vec{AB} + \vec{AD} + \vec{AA'})$.}
	\option{$\vec{AO} = \frac{1}{4}(\vec{AB} + \vec{AD} + \vec{AA'})$.}
	\option{$\vec{AO} = \frac{2}{3}(\vec{AB} + \vec{AD} + \vec{AA'})$.}
\end{mcq}

\begin{mcq}{Cho hình chóp $S.ABCD$ có đáy $ABCD$ là hình bình hành. Đặt $\vec{SA} = \vec{a}, \vec{SB} = \vec{b}, \vec{SC} = \vec{c}, \vec{SD} = \vec{d}$. Mệnh đề nào sau đây là đúng?\image{0.6}{}{lt_6.png}}{0}{1}{Gọi $O$ là tâm hình bình hành $ABCD$, ta có $\vec{SA} + \vec{SC} = 2\vec{SO} = \vec{SB} + \vec{SD} \Rightarrow \vec{a} + \vec{c} = \vec{b} + \vec{d}$.}
	\option{$\vec{a} + \vec{c} = \vec{b} + \vec{d}$.}
	\option{$\vec{a} + \vec{c} + \vec{b} + \vec{d} = \vec{0}$.}
	\option{$\vec{a} + \vec{d} = \vec{b} + \vec{c}$.}
	\option{$\vec{a} + \vec{b} = \vec{c} + \vec{d}$.}
\end{mcq}

\begin{mcq}{Cho hình lập phương $ABCD.A'B'C'D'$. Góc giữa hai vectơ $\vec{A'B}$ và $\vec{CD}$ là:\image{0.6}{}{lt_7.png}}{2}{1}{Ta có $\vec{CD} = \vec{BA}$. Do đó $(\vec{A'B}, \vec{CD}) = (\vec{A'B}, \vec{BA}) = 180^\circ - \widehat{A'BA} = 180^\circ - 45^\circ = 135^\circ$.}
	\option{$45^\circ$.}
	\option{$60^\circ$.}
	\option{$135^\circ$.}
	\option{$120^\circ$.}
\end{mcq}

\begin{mcq}{Cho hình hộp $ABCD.A'B'C'D'$ có $\vec{AB} = \vec{a}, \vec{AC} = \vec{b}, \vec{AA'} = \vec{c}$. Gọi $I$ là trung điểm của $B'C'$, $K$ là giao điểm của $A'I$ và $B'D'$. Mệnh đề nào sau đây là đúng?\image{0.6}{}{lt_8.png}}{0}{1}{Ta có $\vec{DK} = \vec{DD'} + \vec{D'K} = \vec{c} + \frac{2}{3}\vec{D'B'} = \vec{c} + \frac{2}{3}(\vec{AB} - \vec{AD}) = \vec{c} + \frac{2}{3}(2\vec{a} - \vec{b}) = \frac{1}{3}(4\vec{a} - 2\vec{b} + 3\vec{c})$.}
	\option{$\vec{DK} = \frac{1}{3}(4\vec{a} - 2\vec{b} + 3\vec{c})$.}
	\option{$\vec{DK} = -\frac{1}{3}(4\vec{a} - 2\vec{b} + \vec{c})$.}
	\option{$\vec{DK} = 4\vec{a} - 2\vec{b} + \vec{c}$.}
	\option{$\vec{DK} = 4\vec{a} - 2\vec{b} + 3\vec{c}$.}
\end{mcq}

\begin{mcq}{Cho hình hộp $ABCD.A'B'C'D'$. Gọi $M$ là điểm trên cạnh $AC$ sao cho $AC = 3MC$. Lấy điểm $N$ trên đoạn $C'D$ sao cho $C'N = x \cdot C'D$. Với giá trị nào của $x$ thì $MN$ song song $BD'$?\image{0.6}{}{lt_9.png}}{0}{1}{Phân tích theo các vectơ cơ sở ta tìm được $x = \frac{2}{3}$ để $\vec{MN} = \frac{1}{3}\vec{BD'}$.}
	\option{$x = \frac{2}{3}$.}
	\option{$x = \frac{1}{3}$.}
	\option{$x = \frac{1}{4}$.}
	\option{$x = \frac{1}{2}$.}
\end{mcq}

\begin{mcq}{Cho tứ diện đều $ABCD$ có $H$ là trung điểm $AB$. Góc giữa 2 vectơ $\vec{DH}, \vec{AD}$ là:\image{0.6}{}{lt_10.png}}{1}{1}{Tam giác $ABD$ đều nên $\widehat{ADH} = 30^\circ$. Suy ra $(\vec{DH}, \vec{AD}) = 180^\circ - \widehat{ADH} = 150^\circ$.}
	\option{$30^\circ$.}
	\option{$150^\circ$.}
	\option{$90^\circ$.}
	\option{$120^\circ$.}
\end{mcq}

\begin{mcq}{Cho tứ diện $ABCD$ có $AB = AC = AD$ và góc $\widehat{BAC} = \widehat{BAD} = 60^\circ$. Góc giữa 2 vectơ $\vec{AB}, \vec{CD}$ là:\image{0.6}{}{lt_11.png}}{2}{1}{Ta có $\vec{AB} \cdot \vec{CD} = \vec{AB} \cdot (\vec{AD} - \vec{AC}) = AB \cdot AD \cdot \cos 60^\circ - AB \cdot AC \cdot \cos 60^\circ = 0 \Rightarrow (\vec{AB}, \vec{CD}) = 90^\circ$.}
	\option{$30^\circ$.}
	\option{$45^\circ$.}
	\option{$90^\circ$.}
	\option{$120^\circ$.}
\end{mcq}

\begin{mcq}{Cho tứ diện $ABCD$. Mệnh đề nào sau đây là đúng?\image{0.6}{}{lt_12.png}}{1}{1}{Ta có $\vec{AB} \cdot \vec{CD} = (\vec{AC} - \vec{BC}) \cdot \vec{CD} = \vec{AC} \cdot \vec{CD} - \vec{BC} \cdot \vec{CD} = \vec{AC} \cdot \vec{CD} + \vec{BC} \cdot \vec{DC}$.}
	\option{$\vec{AB} \cdot \vec{CD} = \vec{AC} \cdot \vec{CD} + \vec{BC} \cdot \vec{CD}$.}
	\option{$\vec{AB} \cdot \vec{CD} = \vec{AC} \cdot \vec{CD} + \vec{BC} \cdot \vec{DC}$.}
	\option{$\vec{AB} \cdot \vec{CD} = \vec{AC} \cdot \vec{CD} + \vec{CB} \cdot \vec{DC}$.}
	\option{$\vec{AB} \cdot \vec{CD} = \vec{AC} \cdot \vec{DC} + \vec{BC} \cdot \vec{CD}$.}
\end{mcq}

\begin{mcq}{Cho hình hộp $ABCD.A'B'C'D'$. Mệnh đề nào sau đây là sai?\image{0.6}{}{lt_13.png}}{1}{1}{Theo quy tắc hình hộp, $\vec{AB} + \vec{A'D'} + \vec{AA'} = \vec{AB} + \vec{AD} + \vec{AA'} = \vec{AC'} \ne \vec{0}$.}
	\option{$\vec{AB} + \vec{DD'} + \vec{C'D'} = \vec{CC'}$.}
	\option{$\vec{AB} + \vec{A'D'} + \vec{AA'} = \vec{0}$.}
	\option{$\vec{AB} + \vec{CD'} - \vec{CC'} = \vec{0}$.}
	\option{$\vec{BC} - \vec{CC'} + \vec{D'C'} = \vec{A'C}$.}
\end{mcq}

\begin{mcq}{Cho hình lập phương $ABCD.A'B'C'D'$ có cạnh bằng $a$. Giá trị của $(\vec{AB} + \vec{AD}) \cdot \vec{AA'}$ là:\image{0.6}{}{lt_14.png}}{0}{1}{Ta có $(\vec{AB} + \vec{AD}) \cdot \vec{AA'} = \vec{AC} \cdot \vec{AA'} = 0$ vì $AA' \perp (ABCD)$.}
	\option{$0$.}
	\option{$a$.}
	\option{$a^2$.}
	\option{$a^2\sqrt{2}$.}
\end{mcq}

\begin{mcq}{Cho tứ diện $ABCD$ có $AC, BD$ cùng vuông góc với $AB$. Gọi $M, N$ lần lượt là trung điểm của hai cạnh $AB, CD$. Mệnh đề nào sau đây là đúng?\image{0.6}{}{lt_15.png}}{2}{1}{Ta có $2\vec{MN} = (\vec{MA} + \vec{MB}) + (\vec{AC} + \vec{BD}) + (\vec{CN} + \vec{DN}) = \vec{AC} + \vec{BD} \Rightarrow \vec{MN} = \frac{1}{2}(\vec{AC} + \vec{BD})$.}
	\option{$\vec{MN} = \vec{AC} + \vec{BD}$.}
	\option{$\vec{MN} = \vec{AD} + \vec{BC}$.}
	\option{$\vec{MN} = \frac{1}{2}(\vec{AC} + \vec{BD})$.}
	\option{$\vec{MN} = \frac{1}{2}(\vec{CA} + \vec{BD})$.}
\end{mcq}

\begin{mcq}{Cho $|\vec{a}| = 4, |\vec{b}| = 5$ và $(\vec{a}, \vec{b}) = 60^\circ$. Giá trị của $|\vec{a} - \vec{b}|^2$ là:}{1}{1}{Ta có $|\vec{a} - \vec{b}|^2 = |\vec{a}|^2 + |\vec{b}|^2 - 2|\vec{a}||\vec{b}|\cos 60^\circ = 16 + 25 - 2 \cdot 4 \cdot 5 \cdot \frac{1}{2} = 21$.}
	\option{$41 - 20\sqrt{3}$.}
	\option{$21$.}
	\option{$61$.}
	\option{$41 + 20\sqrt{3}$.}
\end{mcq}

\begin{mcq}{Cho hình lập phương $ABCD.A'B'C'D'$ có cạnh bằng $a$. Giá trị của $\vec{AB'} \cdot \vec{BD}$ là:\image{0.6}{}{lt_16.png}}{3}{1}{Ta có $\vec{AB'} \cdot \vec{BD} = (\vec{AB} + \vec{AA'}) \cdot (\vec{AD} - \vec{AB}) = -\vec{AB}^2 = -a^2$.}
	\option{$a^2$.}
	\option{$-a^2\sqrt{3}$.}
	\option{$0$.}
	\option{$-a^2$.}
\end{mcq}

\begin{mcq}{Cho hình lập phương $ABCD.A'B'C'D'$ có cạnh bằng $a$. Giá trị của $(\vec{AB} + \vec{AD}) \cdot \vec{AB'}$ là:\image{0.6}{}{lt_17.png}}{2}{1}{Ta có $(\vec{AB} + \vec{AD}) \cdot \vec{AB'} = \vec{AC} \cdot (\vec{AB} + \vec{AA'}) = \vec{AC} \cdot \vec{AB} = a\sqrt{2} \cdot a \cdot \cos 45^\circ = a^2$.}
	\option{$0$.}
	\option{$a$.}
	\option{$a^2$.}
	\option{$a^2\sqrt{2}$.}
\end{mcq}

\begin{mcq}{Gọi $AD$ là phân giác của $\widehat{BAC}$. Biết $\vec{AB} \cdot \vec{AD} = 2, \vec{AC} \cdot \vec{AD} = 4, |\vec{AB}| = 1$, khi đó giá trị của $|\vec{AC}|$ là:\image{0.6}{}{lt_18.png}}{1}{1}{Đặt góc phân giác là $\alpha$, ta có $AD \cos\alpha = 2$ và $|\vec{AC}| \cdot AD \cos\alpha = 4 \Rightarrow |\vec{AC}| = 2$.}
	\option{$1$.}
	\option{$2$.}
	\option{$3$.}
	\option{$4$.}
\end{mcq}

\begin{mcq}{Cho tứ diện $SABC$ có $SA, SB, SC$ đôi một vuông góc và $SA = SB = SC = a$. Gọi $M$ là trung điểm của $BC$. Góc giữa hai vectơ $\vec{SM}$ và $\vec{AB}$ là:\image{0.6}{}{lt_19.png}}{3}{1}{Ta có $\vec{SM} \cdot \vec{AB} = \frac{1}{2}(\vec{SB} + \vec{SC}) \cdot (\vec{SB} - \vec{SA}) = \frac{a^2}{2}$. Suy ra $\cos(\vec{SM}, \vec{AB}) = \frac{a^2/2}{(a\sqrt{2}/2)(a\sqrt{2})} = \frac{1}{2} \Rightarrow (\vec{SM}, \vec{AB}) = 60^\circ$.}
	\option{$30^\circ$.}
	\option{$45^\circ$.}
	\option{$90^\circ$.}
	\option{$60^\circ$.}
\end{mcq}

\end{quiz}

% =========================================================================
% 2. CÂU TRẮC NGHIỆM ĐÚNG SAI
% =========================================================================

\begin{quiz}{Câu trắc nghiệm đúng sai}

\begin{truefalse}{Cho hình hộp $ABCD.A'B'C'D'$. Xét tính đúng sai của các khẳng định sau:\image{0.6}{}{lt_20.png}}{1}{}
	\statement{false}{$\vec{AA'} + \vec{D'C'} - \vec{BC} = \vec{AC'}$.}
	\statement{false}{$\vec{AB} = \vec{A'B'} = \vec{CD} = \vec{D'C'}$.}
	\statement{false}{$\vec{CC'} + \vec{C'D'} + \vec{BC} = \vec{AC'}$.}
	\statement{true}{$\vec{AA'} + \vec{D'C'} + \vec{BC} = \vec{AC'}$.}
\end{truefalse}

\begin{truefalse}{Cho hình lăng trụ $ABC.A'B'C'$. Gọi $I$ là trung điểm $B'C'$ và $G'$ là trọng tâm tam giác $A'B'C'$. Xét tính đúng sai của các khẳng định sau:\image{0.6}{}{lt_21.png}}{1}{}
	\statement{false}{$\vec{AG'} = \vec{AA'} + \frac{1}{3}\vec{AB} + \frac{2}{3}\vec{AC}$.}
	\statement{true}{$\vec{AG'} = \vec{AA'} + \frac{1}{3}\vec{AB} + \frac{1}{3}\vec{AC}$.}
	\statement{true}{$\vec{A'I} = -\frac{1}{2}\vec{AA'} + \frac{1}{2}\vec{AB'} + \frac{1}{2}\vec{AC}$.}
	\statement{false}{$\vec{A'I} = \frac{1}{2}\vec{AC'} + \frac{1}{2}\vec{A'B}$.}
\end{truefalse}

\begin{truefalse}{Cho hình lập phương $ABCD.A'B'C'D'$ có cạnh bằng $a$. Xét tính đúng sai của các khẳng định sau:\image{0.6}{}{lt_22.png}}{1}{}
	\statement{false}{$(\vec{AC'}, \vec{AC}) = 45^\circ$.}
	\statement{true}{$\vec{BB'} \cdot \vec{AD} = 0$.}
	\statement{true}{$(\vec{AB} + \vec{AD}) \cdot \vec{B'D'} = 0$.}
	\statement{false}{$\vec{CC'} + \vec{C'D'} + \vec{BC} = \vec{AC'}$.}
\end{truefalse}

\begin{truefalse}{Cho hình lập phương $ABCD.A'B'C'D'$ có cạnh bằng $a$. Xét tính đúng sai của các khẳng định sau:\image{0.6}{}{lt_23.png}}{1}{}
	\statement{false}{$\vec{AC'} + \vec{CA'} + 2\vec{CC'} = \vec{0}$.}
	\statement{false}{Độ dài vectơ $\vec{x} = \vec{AA'} + \vec{AC'}$ theo $a$ bằng $\frac{a\sqrt{6}}{2}$.}
	\statement{true}{$(\vec{BD}, \vec{B'C}) = 60^\circ$.}
	\statement{true}{$\vec{BC} - \vec{CC'} + \vec{D'C'} = \vec{A'C}$.}
\end{truefalse}

\begin{truefalse}{Cho tứ diện $ABCD$. Xét tính đúng sai của các khẳng định sau:\image{0.6}{}{lt_24.png}}{1}{}
	\statement{false}{$\vec{AB} \cdot \vec{CD} = \vec{AC} \cdot \vec{CD} + \vec{CB} \cdot \vec{DC}$.}
	\statement{false}{$\vec{AB} + \vec{CD} = \vec{AC} + \vec{DB}$.}
	\statement{true}{$(\vec{AB}, \vec{CD}) = 90^\circ$ nếu $ABCD$ là tứ diện đều.}
	\statement{true}{Giả sử tứ diện $ABCD$ có $AC, BD$ cùng vuông góc với $AB$. Gọi $M, N$ lần lượt là trung điểm của hai cạnh $AB, CD$. Khi đó $\vec{MN} = \frac{1}{2}(\vec{AC} + \vec{BD})$.}
\end{truefalse}

\end{quiz}

% =========================================================================
% 3. CÂU TRẮC NGHIỆM TRẢ LỜI NGẮN
% =========================================================================

\begin{quiz}{Câu trắc nghiệm trả lời ngắn}

\shortanswer{Cho hình chóp tứ giác đều $S.ABCD$ có độ dài tất cả các cạnh bằng $1$.\image{0.6}{}{lt_25.png}

a) Tính tích vô hướng $\vec{AS} \cdot \vec{BC}$.}{$0{,}5$}{1}{Ta có $\vec{AS} \cdot \vec{BC} = \vec{AS} \cdot \vec{AD} = 1 \cdot 1 \cdot \cos 60^\circ = 0{,}5$.}

\shortanswer{Cho hình chóp tứ giác đều $S.ABCD$ có độ dài tất cả các cạnh bằng $1$.

b) Tính tích vô hướng $\vec{AS} \cdot \vec{AC}$.}{$1$}{1}{Tam giác $SAC$ vuông cân tại $S$ nên $\widehat{SAC} = 45^\circ$. Do đó $\vec{AS} \cdot \vec{AC} = 1 \cdot \sqrt{2} \cdot \cos 45^\circ = 1$.}

\shortanswer{Cho hình lập phương $ABCD.EFGH$ có cạnh bằng $2$.\image{0.6}{}{lt_26.png}

a) Tính $\vec{BC} \cdot \vec{AH}$.}{$4$}{1}{Ta có $\vec{BC} \cdot \vec{AH} = \vec{AD} \cdot \vec{AH} = 2 \cdot 2\sqrt{2} \cdot \cos 45^\circ = 4$.}

\shortanswer{Cho hình lập phương $ABCD.EFGH$ có cạnh bằng $2$.

b) Tính $\vec{AF} \cdot \vec{EG}$.}{$4$}{1}{Tam giác $AFC$ đều nên $\vec{AF} \cdot \vec{EG} = \vec{AF} \cdot \vec{AC} = 2\sqrt{2} \cdot 2\sqrt{2} \cdot \cos 60^\circ = 4$.}

\shortanswer{Cho hình lập phương $ABCD.EFGH$ có cạnh bằng $2$.

c) Tính $\vec{AC} \cdot \vec{FE}$.}{$-4$}{1}{Ta có $\vec{AC} \cdot \vec{FE} = -\vec{AC} \cdot \vec{AB} = -(2\sqrt{2}) \cdot 2 \cdot \cos 45^\circ = -4$.}

\shortanswer{Nếu một vật có khối lượng $m$ ($\text{kg}$) thì lực hấp dẫn $\vec{P}$ của Trái Đất tác dụng lên vật được xác định theo công thức $\vec{P} = m\vec{g}$, trong đó $\vec{g}$ là gia tốc rơi tự do có độ lớn $g = 9{,}8\text{ m/s}^2$. Tính độ lớn của lực hấp dẫn của Trái Đất tác dụng lên một quả táo có khối lượng $102\text{ gam}$ (làm tròn đến hàng đơn vị, đơn vị: N).}{$1$}{1}{Đổi $102\text{ g} = 0{,}102\text{ kg}$. Ta có $P = m \cdot g = 0{,}102 \cdot 9{,}8 = 0{,}9996\text{ N} \approx 1\text{ N}$.}

\shortanswer{Trọng lực $\vec{P}$ của Trái Đất tác dụng lên vật được tính bởi công thức $\vec{P} = m\vec{g}$, trong đó $m$ là khối lượng của vật (đơn vị: $\text{kg}$); $\vec{g}$ là gia tốc rơi tự do, có hướng đi xuống và có độ lớn $g = 9{,}8\text{ m/s}^2$. Tính độ lớn của trọng lực tác dụng lên một quả bóng có khối lượng $450\text{ gam}$ (làm tròn đến một chữ số thập phân, đơn vị: N).}{$4{,}4$}{1}{Đổi $450\text{ g} = 0{,}45\text{ kg}$. Ta có $P = m \cdot g = 0{,}45 \cdot 9{,}8 = 4{,}41\text{ N} \approx 4{,}4\text{ N}$.}

\shortanswer{Một chất điểm $A$ nằm trên mặt phẳng nằm ngang $(\alpha)$, chịu tác động bởi ba lực $\vec{F_1}, \vec{F_2}, \vec{F_3}$. Các lực $\vec{F_1}, \vec{F_2}$ có giá nằm trong $(\alpha)$ và $(\vec{F_1}, \vec{F_2}) = 135^\circ$, còn lực $\vec{F_3}$ có giá vuông góc với $(\alpha)$ và hướng lên trên. Xác định độ lớn hợp lực của các lực $\vec{F_1}, \vec{F_2}, \vec{F_3}$ (làm tròn đến một chữ số thập phân, đơn vị: N), biết rằng độ lớn của ba lực đó lần lượt là $20\text{ N}, 15\text{ N}$ và $10\text{ N}$.\image{0.6}{}{lt_27.png}}{$17{,}3$}{1}{Hợp lực $\vec{F} = \vec{F_1} + \vec{F_2} + \vec{F_3}$. Ta có $|\vec{F}|^2 = |\vec{F_1}|^2 + |\vec{F_2}|^2 + 2|\vec{F_1}||\vec{F_2}|\cos 135^\circ + |\vec{F_3}|^2 = 725 - 300\sqrt{2} \approx 300{,}74 \Rightarrow |\vec{F}| \approx 17{,}3\text{ N}$.}

\shortanswer{Người ta treo một vật trang trí $O$ có khối lượng $m = 2\text{ kg}$ trên trần nhà bằng các sợi dây nhẹ, không co dãn tại các điểm $A, B, C$. Để đảm bảo lực phân phối đều trên các dây và tính thẩm mĩ, người ta chọn độ dài các dây sao cho $OABC$ là tứ diện đều. Gọi $\vec{T_1}, \vec{T_2}, \vec{T_3}$ lần lượt là lực căng dây của ba dây treo tại $A, B, C$. Lấy $g = 10\text{ m/s}^2$.\image{0.6}{}{lt_28.png}

a) Tính cường độ của hợp lực các lực căng dây (đơn vị: N).}{$20$}{1}{Vật ở vị trí cân bằng nên hợp lực các lực căng dây cân bằng với trọng lực: $|\vec{T}| = P = m \cdot g = 2 \cdot 10 = 20\text{ N}$.}

\shortanswer{Người ta treo một vật trang trí $O$ có khối lượng $m = 2\text{ kg}$ trên trần nhà bằng các sợi dây nhẹ, không co dãn tại các điểm $A, B, C$. Để đảm bảo lực phân phối đều trên các dây và tính thẩm mĩ, người ta chọn độ dài các dây sao cho $OABC$ là tứ diện đều. Gọi $\vec{T_1}, \vec{T_2}, \vec{T_3}$ lần lượt là lực căng dây của ba dây treo tại $A, B, C$. Lấy $g = 10\text{ m/s}^2$.

b) Tính cường độ của lực căng trên mỗi dây (làm tròn đến một chữ số thập phân, đơn vị: N).}{$8{,}2$}{1}{Ta có $|\vec{T_1} + \vec{T_2} + \vec{T_3}|^2 = 6T^2 = 20^2 = 400 \Rightarrow T = \frac{10\sqrt{6}}{3} \approx 8{,}2\text{ N}$.}

\end{quiz}

\end{document}
